% RAID-Schemata (RAID 0, 1, 4, 5, 6) mit einheitlicher Notation:
% Buchstabe = Streifen, Zahl = Position im Streifen, p/q = Paritätsblock.
% Ersetzt die früheren Wikimedia-Grafiken (Cburnett), damit Skript, Folien
% und Aufgaben dieselben Namen und dieselbe Festplattennummerierung (ab 1) haben.
%
% Aufruf: \raidschema{<Festplatten>}{<Streifen>}{<Zeilen>}
%   <Zeilen>: eine Gruppe pro Streifen, darin pro Festplatte <Typ>/<Name>
%   Typ d = Datenblock, p = Paritätsblock P, q = Paritätsblock Q (RAID 6)
%   Beispiel RAID 4 mit drei Festplatten:
%   \raidschema{3}{2}{{d/A_1, d/A_2, p/A_p}, {d/B_1, d/B_2, p/B_p}}
%
% \providecommand, damit die Datei in Skript und Folien mehrfach
% eingebunden werden darf.
\providecommand{\raidschema}[3]{%
    \begin{tikzpicture}[
            raidblock/.style={draw, minimum width=1.2cm, minimum height=0.55cm, inner sep=1pt, font=\small},
            raidd/.style={raidblock, fill=blue!8},
            raidp/.style={raidblock, fill=orange!35},
            raidq/.style={raidblock, fill=red!25},
            raiddisk/.style={draw, thick, rounded corners=4pt, fill=gray!10}
        ]
        % Festplatten als Säulen, darunter die Nummer
        \foreach \j in {1,...,#1} {
                \draw[raiddisk] ({(\j-1)*1.9-0.8}, 0.5) rectangle ({(\j-1)*1.9+0.8}, {-(#2-1)*0.75-0.5});
                \node[font=\footnotesize, anchor=north] at ({(\j-1)*1.9}, {-(#2-1)*0.75-0.6}) {Festplatte \j};
            }
        % Blöcke: Zeile = Streifen, Spalte = Festplatte
        \foreach \row [count=\i] in {#3} {
                \foreach \typ/\lbl [count=\j] in \row {
                        \node[raid\typ] at ({(\j-1)*1.9}, {-(\i-1)*0.75}) {$\lbl$};
                    }
            }
    \end{tikzpicture}%
}
